
\documentclass[10pt,a4paper,english]{report}
\usepackage{amssymb}




%%%%%%%%%%%%%%%%%%%%%%%%%%%%%%%%%%%%%%%%%%%%%%%%%%%%%%%%%%%%%%%%%%%%%%%%%%%%%%%%%%%%%%%%%%%%%%%%%%%%
%\usepackage{makeidx}
%\usepackage[T1]{fontenc}
%%\usepackage[latin9]{inputenc}
\usepackage{amsmath}
\usepackage{graphicx}
\usepackage{graphics}
\usepackage[dvips]{epsfig}
\usepackage{esint}
\usepackage[authoryear]{natbib}
\usepackage{rotating}
\usepackage{relsize}
\usepackage{Sweave}
\usepackage{xcolor}
\usepackage{babel}
\usepackage{hyperref}



\begin{document}
\title{Computer Exercise 3: Linear Mixed Models}
\author{Lars Ronnegard}
\maketitle
\section*{Introduction}
The data set estrone.data contains 80 observations of estrone levels ($y$) from 5 females (from the book: Pawitan 2001 In all likelihood).


\begin{Schunk}
\begin{Sinput}
y <- c(13.61728, 13.61728, 13.42423, 13.01030, 13.97940, 13.42423, 14.31364, 
    13.97940, 13.42423, 13.42423, 13.61728, 13.61728, 14.31364, 12.78754, 
    13.61728, 12.55273, 13.97940, 15.18514, 14.31364, 14.31364, 14.77121, 
    14.47158, 13.80211, 13.42423, 14.14973, 14.77121, 14.77121, 14.62398, 
    14.62398, 15.68202, 13.80211, 14.47158, 15.79784, 15.79784, 16.12784, 
    15.79784, 15.79784, 15.05150, 15.79784, 16.23249, 15.44068, 16.02060, 
    16.12784, 15.68202, 14.47158, 15.56303, 14.77121, 15.68202, 11.46128, 
    12.04120, 11.76091, 12.78754, 13.01030, 13.42423, 12.04120, 12.78754, 
    12.30449, 12.55273, 13.01030, 12.55273, 10.79181, 12.30449, 11.76091, 
    11.13943, 16.62758, 15.56303, 14.77121, 14.62398, 15.56303, 14.91362, 
    14.77121, 15.05150, 15.05150, 14.91362, 14.77121, 15.05150, 13.97940, 
    14.62398, 14.91362, 15.05150)

id <- c( 1, 1, 1, 1, 1, 1, 1, 1, 1, 1, 1, 1, 1, 1, 1, 1, 
   2, 2, 2, 2, 2, 2, 2, 2, 2, 2,  2, 2, 2, 2, 2, 2, 
   3, 3, 3, 3, 3, 3, 3, 3, 3, 3, 3, 3, 3, 3, 3, 3, 
   4, 4, 4, 4, 4, 4, 4,  4, 4, 4, 4, 4, 4, 4, 4, 4, 
   5, 5, 5, 5, 5, 5, 5, 5, 5, 5, 5, 5, 5, 5, 5, 5 )

#Collect all the data in a data frame
estrone.data <- data.frame(y=y, id=as.factor(id))
\end{Sinput}
\end{Schunk}

\section*{Exercise 1}
Fit a linear mixed model using the lmer() function in the lme4 package with a fixed intercept and female as individual. What is the estimated residual variance?  What is the estimated random effect variance?

\section*{Exercise 2}
Fit the model in Exercise 1 using the hglm2 function in the hglm package. Do you get the same estimates using the hglm2() function and the lmer() function?
\newline
Plot the diagnostic plots.

\section*{Exercise 3}
Construct the design matrix X and the incidence matrix Z in R.
\newline
Fit the model in Exercise 2 using the hglm function in the hglm package. Do you still get the same estimates?
\newline
Plot the diagnostic plots.

\section*{Exercise 4 (If you get time)}
Compute a 95\% confidence interval for the random effect variance.


\end{document}


